% =============================================================================
% PARKING LOT — Content not yet reviewed with Dan
% This file is NOT included in main.tex. It preserves drafted content
% for future use once discussed with the advisor.
%
% Contents:
%   1. Causal identification assumptions (exchangeability, consistency, positivity)
%   2. Observed housing status (direct/partial observations, H^obs, Y^obs)
%   3. Follow-up strata (C = 0,...,5)
%
% NOTE: Some notation here uses the OLD conventions (ordinal H(t), 90-day
% follow-up, A(30), M^v(t)). Will need updating before reuse.
% =============================================================================

\subsection{Assumptions}
\subsubsection{Causal identification}

We make the following standard causal inference assumptions.

\paragraph{Conditional exchangeability.}  First, we assume that, conditional on baseline covariates and baseline housing status, potential outcomes are independent of treatment initation during the grace period:
\begin{align*}
    Y^1(s), Y^0(s) \perp A(s) \mid H(0), \boldsymbol{L}, \boldsymbol{X}(0) \quad \text{for all } s \in [0, 30]
\end{align*}



\paragraph{Consistency.}
\begin{align*}
    Y^a(s) &= Y \quad \text{if } A=a \text{ and } T^A=s
\end{align*}

\paragraph{Positivity.}
\begin{align*}
    0 < P(A(s)=1 \mid H(0), \boldsymbol{L}, \boldsymbol{X}(0)) < 1 \quad \text{for all } H(0), \boldsymbol{L}, \boldsymbol{X}(0) \text{ and } s \in [0, 30]
\end{align*}


\subsubsection{Observed housing status}
Housing status may be fully, partially, or not observed at time $t$ depending on the patient's vital status, program enrollment status, and healthcare visits at time $t$. By fully observed we mean that we can determine the exact value of $H(t)$. By partially observed, we mean that we can determine whether the patient's housing status falls into a certain subset of the four categories, but we cannot determine the exact value of $H(t)$. By not observed, we mean that we cannot directly determine anything about $H(t)$.

\paragraph{Direct observations of housing.} We assume that housing status is fully observed in four settings. First, for all patients at $t=0$ (time of Call Center contact) through the Center's assessment. Second, at the time of death, since $H(t)$ is defined to be $0$ if $D(t)=1$. Third, when a patient is enrolled in GPD, since patients in GPD are considered to be in sheltered homelessness. Fourth, when a patient enters SSVF.

\paragraph{Partial observations of housing.}
Housing status is partially observed for patients living in HUD-VASH housing. For these patients, they are known to be housed, but we cannot determine their risk status.

For patients enrolled in HUD-VASH who are in the process of obtaining HUD-VASH housing, we can determine that they are not stably housed, but we cannot determine whether they are housed with risk of homelessness or homeless during the housing search.

For patients who enrolled in SSVF and exited SSVF prior to follow-up, we can determine whether they exit to unsheltered homelessness, sheltered homelessness, or independent housing, but for the third group we cannot determine whether they are housed with risk of homelessness or stably housed.

Note that stable housing is never directly observed for any patient, since we cannot determine whether a patient is housed with risk of homelessness or stably housed. We can only determine whether a patient is not stably housed (i.e.\ unsheltered homeless, sheltered homeless, or housed with risk of homelessness) or is potentially stably housed (i.e.\ living in HUD-VASH housing or exited SSVF to independent housing).

Define indicators $H^{\text{obs}}(t)$ for whether housing status is observed at time $t$ and $Y^{\text{obs}}$ for whether the primary outcome $Y$ is observed at follow-up.
The observation process for $Y$ is less restrictive than housing since we can determine whether $Y=0$ if we know the patient is deceased or not stably housed, even if we don't know their exact housing status.


% =============================================================================
% Follow-up strata
% =============================================================================

\subsection{Follow-up strata}

Partition patients into strata defined by their vital status and program enrollment at follow-up time $\tau$. Let $C = C(\boldsymbol{L}, \overline{\boldsymbol{X}}(\tau), D(\tau))$ denote stratum membership:

\begin{itemize}
    \item $C = 0$: Deceased ($D(\tau) = 1$)
    \item $C = 1$: Enrolled in GPD ($D(\tau) = 0$, $X_3(\tau) = 1$)
    \item $C = 2$: Enrolled in HUD-VASH but not living in HUD-VASH housing, not enrolled in GPD ($D(\tau) = 0$, $X_1(\tau) = 1$, $X_2(\tau) = 0$, $X_3(\tau) = 0$)
    \item $C = 3$: Living in HUD-VASH housing ($D(\tau) = 0$, $X_2(\tau) = 1$)
    \item $C = 4$: Not enrolled in any VA homeless program at $\tau$, but enrolled in and exited a VA homeless program between $t = 0$ and $\tau$
    \item $C = 5$: Never enrolled in any VA homeless program between $t = 0$ and $\tau$
\end{itemize}

For $C \in \{0, 1, 2\}$, the primary outcome is determined directly: $Y = 0$. For $C \in \{3, 4, 5\}$, $Y$ must be estimated from manifest variables via the latent variable model.
